\chapter{Considera\c{c}\~oes Finais}
\label{cap:consideracoes-finais}

Esta qualifica\c{c}\~ao reuniu o desenvolvimento de m\'etodos de vis\~ao computacional para identificar, localizar e organizar parcelas experimentais de forrageiras em imagens orbitais e de drone e definiu a quest\~ao que orienta a continuidade da pesquisa: como a composi\c{c}\~ao do conjunto de treinamento, em particular a sele\c{c}\~ao de imagens por um crit\'erio quantitativo de instabilidade das predi\c{c}\~oes, afeta a generaliza\c{c}\~ao do detector.

Na avalia\c{c}\~ao comparativa sobre 27 imagens de teste, com 7.852 parcelas, os checkpoints de caixas orientadas obtiveram os maiores valores de mAP@0,5:0,95, entre 0,599 e 0,628. O checkpoint 260525\_2231, indicado para infer\^encia no sistema, alcan\c{c}ou precis\~ao global de 0,952, recall de 0,930 e F1 de 0,941, com F1 de 0,949 nas imagens orbitais e de 0,939 nas de drone. Os outros dois checkpoints de caixas orientadas avaliados nas mesmas condi\c{c}\~oes tiveram recall de 0,523 e 0,546, com a queda concentrada nas imagens de drone. Esses resultados mostram que um detector de caixas orientadas alcan\c{c}ou precis\~ao e recall elevados nesse conjunto, tanto em imagens orbitais quanto de drone, mas n\~ao isolam o efeito da representa\c{c}\~ao, da arquitetura ou dos dados de treinamento, que variaram entre as execu\c{c}\~oes.

Na organiza\c{c}\~ao espacial, o procedimento que combina uma refer\^encia numerada pelo operador, registro geom\'etrico e atribui\c{c}\~ao por custo m\'inimo foi aplicado a 23 imagens de um campo e atribuiu identificador a 90,77\% dos 5.168 objetos. Esse valor mede a completude das atribui\c{c}\~oes; a corre\c{c}\~ao dos identificadores e o tratamento de posi\c{c}\~oes experimentais sem detec\c{c}\~ao permanecem em aberto.

A principal limita\c{c}\~ao dos resultados est\'a nos dados: a independ\^encia entre treinamento e avalia\c{c}\~ao foi reduzida pela divis\~ao aleat\'oria por imagem, os dados exatos de treinamento de parte dos checkpoints s\~ao desconhecidos e n\~ao houve execu\c{c}\~oes repetidas. As m\'etricas descrevem, portanto, o desempenho na parti\c{c}\~ao de teste utilizada e n\~ao estimam a generaliza\c{c}\~ao para campos independentes, restri\c{c}\~ao que alcan\c{c}a tamb\'em o classificador auxiliar, com acur\'acia de 76,66\%.

Essas limita\c{c}\~oes definem as condi\c{c}\~oes da etapa seguinte. O estudo sobre a composi\c{c}\~ao do conjunto de treinamento ser\'a conduzido com o acervo ampliado, com um protocolo de divis\~ao das imagens por grupos espaciais e com as demais condi\c{c}\~oes de treinamento e avalia\c{c}\~ao mantidas constantes, comparando o conjunto completo, subconjuntos aleat\'orios de mesmo tamanho, o limiar de F1 de 0,97 aplicado a valores obtidos fora da amostra e a sele\c{c}\~ao por instabilidade das predi\c{c}\~oes. O suporte \`a hip\'otese de que essa sele\c{c}\~ao supera subconjuntos aleat\'orios de mesmo tamanho ser\'a avaliado pelas diferen\c{c}as estimadas e por seus intervalos de confian\c{c}a.
